\begin{frame}{유도 템플릿 ①: 거꾸로 대입 (backward induction)}
  \begin{itemize}
    \item \kb{뼈대}: 전이가 ``앞으로만'' 흐르는 MDP에서 벨만방정식
      $V^\pi(s)=R(s,\pi(s))+\gamma V^\pi(s')$를 \textbf{터미널에서
      출발해 거꾸로 한 칸씩 대입}하면 \textbf{연립방정식을 풀지 않아도}
      풀린다.
    \item \kb{왜 동작하나}: 전이 그래프가 \textbf{사이클 없이 한 방향}
      (0$\to$1$\to$2 사슬, 유한 스텝 후 ``끝'' 보장)이면 연립방정식의
      마지막 줄이 더 이상 다른 상태를 참조하지 않음 $\to$ 해가 터미널
      에서 ``퍼져 나간다''. (Ch\,3 3칸 보드게임, Ch\,5 블랙잭, 7.3의
      5$\times$5 격자)
  \end{itemize}
  \vskip0.5em
  \begin{block}{작은 예제: 3-state 사슬, $\gamma=0.9$ (0$\to$1: $-1$,
  1$\to$2: $-1$, 2=터미널)}
    \[
      V(2)=0 \ \text{(터미널)} \qquad
      V(1)=-1+0.9\,V(2)=-1 \qquad
      V(0)=-1+0.9\,V(1)=-1.9
    \]
    \textbf{정확성 확인}(답안에 한 줄): $V=(-1.9,-1,0)$를 벨만방정식
    \textbf{양변에 다시 대입}해 좌우가 일치하는가?
  \end{block}
  \vskip0.4em
  {\footnotesize \textbf{자주 하는 실수} --- (1) 터미널 진입 보상 자체를
  빼먹기, (2) \textbf{대입 방향을 거꾸로} 쓰기($V(1)=-1+0.9V(0)$로).
  대입 방향은 ``다음 상태의 값이 \emph{이미 알려진} 쪽'' = 전이 화살표와
  \textbf{반대}.}
\end{frame}
